\chapter{Export}\label{ch:export}
\begin{goals}
\textbf{In this chapter} you write the part of a file you need to a new file: the filtered, sorted view,
the columns you choose, calculated columns included, at any size.
\end{goals}

Export is how something leaves lamina. lamina never changes the file it opened, so to keep a filtered
subset, a converted \file{.dm} or a table with derived columns, you write a new file.

\section{The export dialog}\label{sec:export-dialog}

\menu{Data>Export…} opens it.

\shot[0.55\linewidth]{export}{The export dialog: the format, the header row, which rows, and which
columns.}

\begin{description}
\item[\ui{Format}] \ui{CSV (comma)}, \ui{TSV (tab)}, \ui{semicolon} or \ui{pipe}. It starts on the open
file's own delimiter when that is tab, and on comma otherwise. A TSV is named \file{.tsv}; the others
\file{.csv}.
\item[\ui{Header row}] write the column names as the first line.
\item[\ui{Rows}] shown when a filter or a sort is on: \ui{Current view (filtered / sorted)} writes the
rows you see, in the order you see them; \ui{All rows} writes the whole file in its own order.
\item[\ui{Columns}] a checkbox per column, calculated columns among them, with a search box and
\ui{all}, \ui{none} and \ui{invert}. Hidden columns start unticked; the count, \code{14 / 16}, says how
many will be written.
\end{description}

\ui{Export} asks where to save the file, then writes it, with the progress in the dialog:
\code{exporting… 42\% (1,204,000 rows)}. When it is done the dialog says \code{✓ 1,204,331 rows →
model.csv} and closes. Closing the dialog during an export cancels it.

\section{What is written}\label{sec:export-what}

\begin{itemize}
\item \textbf{Values as the file has them.} Text as it was read; numbers as the file wrote them for
delimited text, so a decimal-comma file exports its numbers with commas (quoted, in a comma-delimited
export). Number formats and colour scales are display only and aren't applied.
\item \textbf{Calculated columns as values}: the export of a calculated column is an ordinary column of
its results.
\item \textbf{Columns in the file's order}, calculated columns after, whatever order the grid shows them
in.
\item \textbf{Quoting} only where needed: a field containing the delimiter, a double quote or a line
break is put in double quotes, with its quotes doubled.
\end{itemize}

\begin{caution}
A single-precision \file{.dm} currently exports its numbers with float noise in the last digits
(\cref{sec:dm}).
\end{caution}

\section{Any size}\label{sec:export-size}

The export reads the view and writes each row as it goes; it holds nothing whole. A filtered view without
a sort is written in the file's own order, block by block, skipping blocks with no matches, which is fast;
a sorted view is written in its sorted order, reading each row where it is.

That needs a save dialog that writes as it goes, which Chrome and Edge have. Where the browser has none
(Firefox), lamina builds the export in memory and downloads it, and stops with \code{too large for the
download fallback} past about 256~million characters; export from Chrome or Edge for larger files. Inside
the lead-acid shell on Android, an export streams into the phone's \file{Downloads} folder, and lamina
offers to share it from there.

\section{Converting}\label{sec:convert}

Since an export reads any format lamina opens and writes delimited text, it is also a converter:

\begin{itemize}
\item \textbf{\file{.dm} to CSV}: open the \file{.dm}, export. Any size.
\item \textbf{Parquet to CSV}: open the Parquet file, export.
\item \textbf{A sheet to CSV}: open the workbook, pick the sheet, export.
\item \textbf{Semicolon and decimal comma to comma and point}: not directly, since values are written as
the file has them. Make calculated columns of the numeric columns (\code{FE * 1}), whose results are
numbers, and export those.
\end{itemize}

\begin{tryit}
Take the high-grade sulphide out of the sample.
\begin{enumerate}
\item Filter \code{LITO = "SULPHIDE" and Cu\_pct >= 1}, and sort by \col{Cu\_pct} descending.
\item Open \menu{Data>Export…}. \ui{Rows} is \ui{Current view}. Click \ui{none}, then tick \col{X},
\col{Y}, \col{Z}, \col{Cu\_pct} and \col{Au\_gpt}.
\item Click \ui{Export} and save. The file has those five columns, the richest block first.
\end{enumerate}
\end{tryit}
